\subsection{A sufficient first-prime condition and its exact quintic spin term}

\begin{result}{E.208}{Proved}
{One zero first diagonal suffices; the remaining prime condition is not a Chebotarev conclusion}
Let $q=q_0$ be any legal first prime in the original recursion
E.102 with the minimal five-adic shell of E.199. Let
$K=\mathbb Q(\zeta_5)$, choose the prime $\mathfrak q$ of $K$
specified by the local frame, and let $\pi_q$ generate it.
Use $U_q=c_q/q$ and the quintic residue logarithm $\ell_q$
of E.202. If
\[
                         \ell_q(3U_q)=0,                         \tag{E.208.1}
\]
then there is a legal original second prime $r=r_0$, with
its prescribed minimal five-adic shell, for which the
matrix E.202.1 is invertible. Its affine first system
therefore has a unique successful root-choice parameter.
Combining it with E.205 gives an ordinary infinite-support
output of the unchanged original factorial recursion whose
continuous coefficient-five third obstruction vanishes.

For a primary generator
$\pi_q\equiv1\pmod{(1-\zeta_5)^3}$, condition E.208.1 is
exactly
\[
               \ell_q(3)+2\ell_q(\sigma_2\pi_q)=0,
               \qquad\sigma_2(\zeta_5)=\zeta_5^2.               \tag{E.208.2}
\]
The analogous second diagonal is
$\ell_r(-2)+2\ell_r(\sigma_2\pi_r)$.
These formulas isolate a conjugate-prime quintic residue
symbol. They do not prove existence of a legal $q$
satisfying E.208.1, and that condition is only sufficient,
not claimed necessary. The simultaneous off-diagonal
prime prescription below is a fixed-field Chebotarev
argument; no such argument for the varying-prime diagonal
is asserted.
\end{result}

\begin{proof}
We first keep the entire original second-prime selection
while arranging both off-diagonal entries to be nonzero.
The classes $z_q$ of E.202 and $w_q$ of E.206 are already
defined over $G_{\mathbb Q,\{5,q,\infty\}}$, independently
of the as yet unchosen $r$. Both have tame residue one
at $q$, and $w_q$ is zero locally at five.
After restriction to $G_K$ they are nonzero homomorphisms
to the additive group $\mathbb F_5$.
Let $L_z/K,L_w/K$ be their cyclic extensions, adjoining
$K$ throughout. They are Galois over $\mathbb Q$.
Conjugation on their degree-five Galois groups has
characters $\chi^2$ and $\chi^3$, respectively.
This follows either directly from the cocycle identity
or from their Kummer eigenspaces $\chi^{-1}$ and
$\chi^{-2}$ in E.202 and E.206.

Those distinct characters make $L_z,L_w$ linearly
disjoint over $K$. Put $L=L_zL_w$.
Then
\[
 \operatorname{Gal}(L/K)=\mathbb F_5(2)\oplus\mathbb F_5(3).
 \tag{E.208.3}
\]
The maximal abelian subextension of $L/\mathbb Q$ is
exactly $K$: its commutator contains both degree-five
factors, since neither character is trivial, and the
quotient by those factors is the abelian group
$\operatorname{Gal}(K/\mathbb Q)$.

At the original second-prime choice, all old local root
coordinates have been selected and the current-prime
condition at $q$ is imposed. Let $F$ be the full
cyclotomic CRT field, enlarged if necessary to contain
$\mu_{120}$. The chosen automorphism of $F$ is the
identity on $\mu_{120}$; it includes the actual
$r\equiv1-5\pmod{25}$ shell. Since $F/\mathbb Q$
is abelian, E.208.3 gives $L\cap F=K$ and hence
\[
 \operatorname{Gal}(LF/F)
     =\mathbb F_5(2)\oplus\mathbb F_5(3).                         \tag{E.208.4}
\]
The original radical is
\[
                  R=F\bigl((3q^{24})^{1/120}\bigr).
\]
As proved in E.102 and used with enlarged cyclotomic
fields in E.205, its Galois group over $F$ is cyclic
of order sixty. Its unique degree-five subextension
is $F((3q^{24})^{1/5})$; the conjugation character on
that cyclic Galois group is $\chi$.

The intersection $LF\cap R$ is Galois over
$\mathbb Q$. Over $F$ its group is both a quotient
of E.208.4 and a cyclic group of exponent dividing
five. If nontrivial it must be the just-named
degree-five subextension of $R$, of character $\chi$.
This is impossible: a quotient of the two-dimensional
module with characters $\chi^2,\chi^3$ has no
$\chi$ constituent. Concretely, the order-four
cyclotomic action is semisimple over $\mathbb F_5$,
with its three distinct nontrivial eigencharacters.
The statement remains true after adjoining all the
CRT roots of unity because that action still surjects
onto $\operatorname{Gal}(K/\mathbb Q)$.
Thus
\[
                              LF\cap R=F.                       \tag{E.208.5}
\]

The E.102 automorphism on $R$, fixing the chosen
$120$th root and restricting to the prescribed CRT
automorphism on $F$, can now be combined with any
two values of $z_q,w_q$ on $LF$. Choose both values
nonzero. Such lifts on $LF$ exist because $L\cap F=K$
and the CRT automorphism is the identity on $K$.
Equation E.208.5 glues the two choices to an actual
automorphism of $LFR$. Conjugacy preserves the two
nonzero conditions, multiplying their coordinates by
nonzero cyclotomic scalars. Chebotarev, in the exact
form used in E.102, therefore supplies infinitely many
eligible $r>q$, avoiding all ramified primes and the
original finite exclusions. Every old CRT condition,
the current residue at $q$, the minimal shell at five,
and the complete splitting in the original radical
field are retained.

For each such $r$, the reciprocity identities E.206.12
give
\[
      M_{21}=z_q(F_r)\ne0,\qquad
      M_{12}=z_r(F_q)=w_q(F_r)\ne0.                              \tag{E.208.6}
\]
The first diagonal is zero by E.208.1. Hence
$\det M=-M_{12}M_{21}\ne0$, independently of the
uncontrolled second diagonal. E.202 gives exactly
one successful parameter, and E.205 completes the
continuous coefficient-five third gate. No condition
at another coefficient prime is used or concluded.

We finish by deriving the diagonal formula from a
primary reciprocity theorem. Put $\lambda=1-\zeta_5$.
By \citet[Lemma1(a)--(b), pp.~322--323]{AuditGrant1996},
every prime generator has an associate congruent to
one modulo $\lambda^3$, and two such associates
differ by a fifth power of a unit. Take this primary
$\pi=\pi_q$. Its four conjugates are primary and
generate distinct primes because $q\equiv1\pmod5$.
For $a=2,3,4$ put
\[
                     s_a=\ell_q(\sigma_a\pi)\in\mathbb F_5.
\]
The main quintic reciprocity law in
\citet[Theorem (Kummer), p.~324]{AuditGrant1996} applies
to each pair $\pi,\sigma_a\pi$. Galois covariance
of the residue symbol consequently gives
\[
 \left(\frac{\sigma_a\pi}{\pi}\right)_5
   =\left(\frac{\pi}{\sigma_a\pi}\right)_5
   =\sigma_a\left(\frac{\sigma_{a^{-1}}\pi}{\pi}\right)_5,
 \qquad s_a=a s_{a^{-1}}.                                       \tag{E.208.7}
\]
Here $a^{-1}$ is taken modulo five. Thus $s_4=4s_4$
implies $s_4=0$, and $s_3=3s_2$.
The exact norm identity $\prod_{a=1}^4\sigma_a\pi=q$
from E.202 gives
\[
 U_q=\prod_{a=2}^4(\sigma_a\pi)^{a-1},\qquad
 \ell_q(U_q)=s_2+2s_3+3s_4=2s_2.                               \tag{E.208.8}
\]
This proves E.208.2. Replacing the rational factor
$3$ by $-2$ gives the second diagonal formula.
Both are independent of primary associates, and
their zero conditions are independent of the
nonzero scalar choice of common Tate frame.

The symbol in E.208.8 has its numerator and
denominator derived from the same varying prime
generator. The fixed extensions $L_z,L_w$ used
after choosing $q$ do not turn this condition on
$q$ into a prescribed Frobenius in a fixed field.
No distribution theorem for that symbol with the
original first-prime restrictions has been proved
or invoked here. This is the precise remaining
step of this sufficient construction. E.213 subsequently
removes the need for that sufficient premise: two legal
candidate first primes and one common legal second prime
produce an invertible first matrix without determining
either varying-prime spin.
\end{proof}
